\section[L. 净化行动 3]{\myboxxii{      \BF{\TT{L}}} 净化行动 3}

\textit{运行时间限制}   2 seconds

\textit{内存使用限制} 1 Gigabyte

\textit{读入文件} \qquad 标准输入

\textit{输出文件} \qquad 标准输出


\begin{quote}
	\textit{除恶务尽——何为「恶」，我们自有裁定。}
\end{quote}

\subsection{题目描述}

Maddy 仍在清剿 Baddy 的邪恶军团……（这已经是第三次了！）

游戏在一个无限大的平面直角坐标系（可以看作无限大的棋盘）上进行。Baddy 操控了 $n$ 个\BF{炮}棋子。这些炮分别部署在 $(x_1, y_1), (x_2, y_2), \ldots, (x_n, y_n)$ 上，且固定不会移动。

在象棋规则中，炮的攻击规则是"隔山打牛"。如果两个炮 $A$ 和 $B$ 处于同一行或同一列，且它们之间\BF{有且仅有}一个其他炮 $C$，那么我们称 $A$ 攻击了 $B$（同时 $B$ 也攻击了 $A$），而 $C$ 就是它们攻击的"炮架"。

每一次，Maddy 可以选择一个\BF{当前没有受到任何其他炮攻击}的炮，并将其摧毁。被摧毁的炮会被立即移出棋盘，它将不再能攻击其他炮，也不能再作为其他炮的"炮架"。

Maddy 的目标是摧毁所有的 $n$ 个炮。是否存在一种摧毁顺序，使得 Maddy 可以最终清空棋盘？

\subsection{输入格式}

第一行输入一个整数 $T$ ($1 \le T \le 10^5$)，表示测试数据的数量。

每组测试数据第一行输入一个整数 $n$ ($1 \le n \le 2 \times 10^5$)，表示炮的数量。

随后 $n$ 行，每行输入两个整数 $x_i, y_i$ ($-10^9 \le x_i, y_i \le 10^9$)，表示第 $i$ 个炮的位置。

保证所有测试数据的 $\sum n \le 2 \times 10^5$，且单组测试数据中所有炮的位置两两不同。

\subsection{输出格式}

对于每组测试数据，若存在一种合法的摧毁顺序能够清空棋盘，则输出一行 \TT{YES}，否则输出一行 \TT{NO}。

注意评测时不会区分 \TT{YES} 和 \TT{NO} 的大小写，换言之当答案是肯定的时候输出 \TT{yes}、\TT{YES}、\TT{Yes}、\TT{YeS} 等都会被认为是正确的。

\subsection{样例 1}

\begin{minipage}{\linewidth}
	\begin{minipage}[t]{0.49\linewidth}
		\begin{lstlisting}[title=\texttt{standard input}]
3
3
0 0
0 1
0 2
5
0 0
0 1
0 2
1 1
2 1
4
0 0
1 0
2 0
3 0
\end{lstlisting}
	\end{minipage}
	\begin{minipage}[t]{0.49\linewidth}
		\begin{lstlisting}[title=\texttt{standard output}]
YES
YES
NO
\end{lstlisting}
	\end{minipage}
\end{minipage}


\subsection{样例 1 解释}

对于第 $1$ 组测试数据，三个炮在同一列上。初始时，$(0, 0)$ 和 $(0, 2)$ 互相攻击（以 $(0, 1)$ 为炮架），而 $(0, 1)$ 没有受到攻击。因此可以先摧毁 $(0, 1)$，然后 $(0, 0)$ 和 $(0, 2)$ 都不再受到攻击，可以按任意顺序摧毁。

对于第 $3$ 组测试数据，四个炮在同一行上。初始时，$(0, 0)$ 和 $(2, 0)$ 互相攻击（以 $(1, 0)$ 为炮架），$(1, 0)$ 和 $(3, 0)$ 互相攻击（以 $(2, 0)$ 为炮架）。因此 $(0, 0)$、$(1, 0)$、$(2, 0)$、$(3, 0)$ 都受到攻击，无法摧毁任何一个，因此无解。
